% ! TeX root = distributed-systems-final-report.tex
\section{Implementation}\label{implementation}


\todo{notes:}

In addition to what said in \ref{req:distr_transparency}, \textbf{scaling} can
easily be performed thanks to the publish-subscribe model: no additional
configuration would be required for the communication with vehicles, although
increasing the amount of controller nodes would require adding some complexity
to the system. Moreover, the MQTT protocol is specifically designed to work with
constrained devices, to grant netowrk reliability. On the other hand, in its
current implementation, the system isn't expected to expand much: the roundabout
surrounding can only contain a limited amount of cars, and the old ones are
constantly replaced, in the controller's memory, as they drive away.

Using standardized communication protocols (JSON messages over MQTT), the system
is open to potentially welcome external vehicles' implementations, for example
even with different languages and technologies: any vehicle with similar
characteristics may join the network and adhere to the controller's supervision.

TLS authentication, already supported by Mosquitto MQTT, could easily be added
in the future.


\subsection{Technological details}\label{technological-details}

